\section{Related Work}

The framework developed here draws on and extends several established research programs.

In \textbf{organizational theory}, the fission principle recapitulates elements of Mintzberg's \cite{mintzberg1979} structural configurations (simple structure, machine bureaucracy, adhocracy) and the coordination mechanisms catalogued by Malone and Crowston \cite{malone1994}. The principal difference is temporal: in agentic systems, structural reconfiguration can occur within minutes rather than months, enabling a much faster evolutionary cycle.

In \textbf{cybernetics}, Ashby's Law of Requisite Variety \cite{ashby1956} states that a controller must possess at least as much variety as the system it regulates. Each fission event can be understood as increasing the system's regulatory variety by dedicating a channel to a dimension of variation previously handled implicitly. The Monitor role, which compresses raw signals into summary metrics, implements a version of Conant and Ashby's \cite{conant1970} Good Regulator theorem: every good regulator of a system must contain a model of that system.

In \textbf{complex adaptive systems}, Kauffman's \cite{kauffman1993} theory of self-organization at the edge of chaos in NK fitness landscapes predicts that systems with intermediate levels of epistatic coupling evolve most effectively. Our observation that optimal fission granularity depends on context window size, task complexity, and communication bandwidth is consistent with this prediction, though we have not verified the correspondence formally.

In \textbf{information security}, the verification regress has been confronted with particular urgency. Thompson's \cite{thompson1984} Turing Award lecture demonstrated that a compiler can harbor a self-propagating trojan invisible to source-level inspection---an engineering manifestation of the incompleteness results we invoke in Section~6. The Trusted Computing Base (TCB) paradigm \cite{rushby1981} responds with an explicitly foundationalist strategy: minimize the set of components that must be assumed trustworthy, then build verification chains outward (hardware root of trust $\to$ firmware $\to$ bootloader $\to$ kernel). This mirrors our conservative GP prior $\mu_0$. Decentralized alternatives such as PGP's Web of Trust and Byzantine fault-tolerant consensus \cite{lamport1982} adopt a coherentist posture: trust emerges from mutual verification among diverse parties, though collectively compromised toolchains can defeat such checks. Bug-bounty programs and penetration testing represent the pragmatist pole: trust is calibrated not by proof but by the empirical absence of successful attack, much as our Trust GP updates on outcome-based preference pairs. Modern Zero Trust Architecture \cite{rose2020} synthesizes all three strategies---hardware roots, cross-layer consistency checks, and continuous behavioral monitoring---producing a dynamic, context-dependent trust posture structurally analogous to the learned HITL boundary of Section~4.1.  TCB minimization also offers an instructive dual to our fission principle: where fission \emph{distributes} trust across differentiated roles, TCB design seeks to \emph{concentrate} the trust-critical surface into the smallest possible core.

In \textbf{multi-agent systems}, the growing literature on LLM-based multi-agent architectures has explored role assignment, communication protocols, and evaluation frameworks. Our contribution is to situate these engineering choices within a broader evolutionary and epistemological framework, connecting them to foundational results in social choice theory, Bayesian optimization, and mathematical logic.
